% ============================================================
% SECTION 4 - CONDITIONING THE CALIBRATION ON THE SEASON
% ============================================================
\section{Conditioning the Calibration on the Season}
\label{sec:algorithm}

Proposition~\ref{prop:pooled_gap} says that a calibration buffer which
ignores the season carries a bias that data cannot remove.  It does not
say how the season should enter.  This section sets out the two ways of
letting it enter that the literature suggests, shows that they are the
two endpoints of a one-parameter family, and isolates two finite-sample
defects that are easy to mistake for properties of seasonal calibration
but belong to the calibration step of the original algorithm.

The first way is stratification: replace the single buffer by $\Sper$
parallel buffers, one per season, so that each interval is calibrated
exclusively on residuals drawn from the same seasonal distribution as
the prediction target.  This is \texttt{EnbPI-S}, described in
Section~\ref{sec:algorithm:description}.  The second is normalisation:
divide each residual by an estimate of its own season's scale and
calibrate on the pooled buffer of standardised residuals, rescaling the
result by the scale of the target season.  This is the seasonal
counterpart of normalised conformal prediction
\citep{dewolf2025conditional} and is described in
Section~\ref{sec:algorithm:normalised}.  Section~\ref{sec:algorithm:family}
combines them.

%------------------------------------------------------------
\subsection{Algorithm Description}
\label{sec:algorithm:description}
%------------------------------------------------------------

\texttt{EnbPI-S} inherits the training phase of Algorithm~1 in
\citet{xu2023conformal} without change.  All $B$ bootstrap models are
fitted once on the full training set, and all $T$ LOO residuals are
computed as in Section~\ref{sec:model:enbpi}.  The only modification
concerns how residuals are stored and used during prediction.

At the end of the training phase, instead of collecting all $T$ LOO
residuals into a single ordered multiset $\hat{\boldsymbol\epsilon}$,
each residual is assigned to its season-specific buffer:
\[
  \hat{\boldsymbol\epsilon}^{(s)}
  \;=\;
  \bigl\{\hat{\epsilon}_t : t \in \calT_s\bigr\},
  \qquad s = 1, \ldots, \Sper,
\]
where $\calT_s = \{t \in [T]: s(t) = s\}$ is the stratum of
season $s$ (Definition~\ref{def:season}), and each buffer
$\hat{\boldsymbol\epsilon}^{(s)}$ has cardinality
$T_s = |\calT_s|$.

At test time $t > T$ with season $\sstar = s(t)$, the point forecast
is computed exactly as in EnbPI,
\[
  \hat{f}_{-t}^{\,\phi}(x_t)
  \;=\;
  \phi\!\bigl(\hat{f}_{-i}^{\,\phi}(x_t) : i = 1,\ldots,T\bigr),
\]
whereas the line search for the optimal lower quantile level now uses
only the season-$\sstar$ buffer:
\begin{equation}
  \label{eq:betahat_strat}
  \hat{\beta}^{(\sstar)}
  \;=\;
  \operatorname*{arg\,min}_{\beta \in [0,\,\alpha]}
  \Bigl[
    \Quant{1-\alpha+\beta}{\hat{\boldsymbol\epsilon}^{(\sstar)}}
    -
    \Quant{\beta}{\hat{\boldsymbol\epsilon}^{(\sstar)}}
  \Bigr].
\end{equation}
The stratified prediction interval is then:
\begin{equation}
  \label{eq:enbpis_interval}
  \hat{C}_t^{\,\alpha,\textup{strat}}
  \;=\;
  \Bigl[
    \hat{f}_{-t}^{\,\phi}(x_t)
    + \Quant{\hat{\beta}^{(\sstar)}}{\hat{\boldsymbol\epsilon}^{(\sstar)}},\;
    \hat{f}_{-t}^{\,\phi}(x_t)
    + \Quant{1-\alpha+\hat{\beta}^{(\sstar)}}{\hat{\boldsymbol\epsilon}^{(\sstar)}}
  \Bigr].
\end{equation}

Every $s_0 \geq 1$ steps, once observations $\{y_j\}_{j=t-s_0}^{t-1}$
become available, only the buffer corresponding to their season is
updated.  For each revealed observation $y_j$ at season
$s_j = s(j)$:
\begin{equation}
  \label{eq:slide_strat}
  \hat{\boldsymbol\epsilon}^{(s_j)}
  \;\leftarrow\;
  \Bigl(
    \hat{\boldsymbol\epsilon}^{(s_j)}
    \setminus
    \{\hat{\epsilon}_{\text{oldest}}^{(s_j)}\}
  \Bigr)
  \cup
  \{\hat{\epsilon}_j\},
  \qquad
  \hat{\epsilon}_j = y_j - \hat{f}_{-j}^{\,\phi}(x_j).
\end{equation}
Buffers for other seasons are untouched.  After the update, each
buffer maintains its original size $T_{s_j}$, so that the effective
calibration history for each season grows at rate $1/\Sper$ relative
to the passage of time.  The complete procedure is given in
Algorithm~\ref{alg:enbpis}.

%------------------------------------------------------------
\begin{algorithm}[t]
\caption{Ensemble Batch Prediction Intervals, Stratified
         (\texttt{EnbPI-S})}
\label{alg:enbpis}
\DontPrintSemicolon
\SetKwInOut{Require}{Require}
\SetKwInOut{Ensure}{Ensure}
\SetKwComment{Comment}{$\triangleright$\ }{}

\Require{Training data $\{(x_i, y_i)\}_{i=1}^T$; prediction
  algorithm $\calA$; significance level $\alpha$; aggregation
  function $\phi$; number of bootstrap models $B$; batch size
  $s_0 \geq 1$; seasonal period $\Sper$; test data
  $\{(x_t, y_t)\}_{t=T+1}^{T+T_1}$
  ($y_t$ revealed only after prediction at $t$).}

\Ensure{Stratified prediction intervals
  $\{\hat{C}_t^{\alpha,\textup{strat}}(x_t)\}_{t=T+1}^{T+T_1}$.}

\BlankLine
\tcp{\textbf{TRAINING PHASE} (identical to \citet[Alg.~1]{xu2023conformal})}
\For{$b = 1, \ldots, B$}{
  Sample index set $S_b = (i_1, \ldots, i_T)$ from $[T]$ with
  replacement\;
  Compute $\hat{f}^{\,b} = \calA\!\bigl((x_i, y_i),\; i \in S_b\bigr)$\;
}
Initialise $\Sper$ empty ordered sets
$\hat{\boldsymbol\epsilon}^{(1)}, \ldots, \hat{\boldsymbol\epsilon}^{(\Sper)}$\;
\For{$i = 1, \ldots, T$}{
  $\hat{f}_{-i}^{\,\phi}(x_i) \leftarrow
    \phi\!\bigl(\hat{f}^{\,b}(x_i):\, b \text{ s.t. } i \notin S_b\bigr)$\;
  $\hat{\epsilon}_i \leftarrow y_i - \hat{f}_{-i}^{\,\phi}(x_i)$\;
  $\hat{\boldsymbol\epsilon}^{(s(i))} \leftarrow
    \hat{\boldsymbol\epsilon}^{(s(i))} \cup \{\hat{\epsilon}_i\}$\;
  \Comment{route residual to its season's buffer}
}

\BlankLine
\tcp{\textbf{PREDICTION PHASE}}
\For{$t = T+1, \ldots, T+T_1$}{
  $\sstar \leftarrow s(t)$\;
  $\hat{f}_{-t}^{\,\phi}(x_t) \leftarrow
    \phi\!\bigl(\hat{f}_{-i}^{\,\phi}(x_t),\; i = 1, \ldots, T\bigr)$\;
  Compute $\hat{\beta}^{(\sstar)}$ as
  $\arg\min_{\beta \in [0,\alpha]}
  \bigl[\Quant{1-\alpha+\beta}{\hat{\boldsymbol\epsilon}^{(\sstar)}}
  -\Quant{\beta}{\hat{\boldsymbol\epsilon}^{(\sstar)}}\bigr]$
  via line search\;
  $w_{\text{lo}} \leftarrow
    \Quant{\hat{\beta}^{(\sstar)}}{\hat{\boldsymbol\epsilon}^{(\sstar)}}$\;
  $w_{\text{hi}} \leftarrow
    \Quant{1-\alpha+\hat{\beta}^{(\sstar)}}{\hat{\boldsymbol\epsilon}^{(\sstar)}}$\;
  Output $\hat{C}_t^{\alpha,\textup{strat}}(x_t) \leftarrow
    [\hat{f}_{-t}^{\,\phi}(x_t) + w_{\text{lo}},\;
     \hat{f}_{-t}^{\,\phi}(x_t) + w_{\text{hi}}]$\;
  \If{$(t - T) \equiv 0 \pmod{s_0}$}{
    \For{$j = t - s_0, \ldots, t - 1$}{
      $\hat{\epsilon}_j \leftarrow y_j - \hat{f}_{-j}^{\,\phi}(x_j)$\;
      Remove oldest element from
        $\hat{\boldsymbol\epsilon}^{(s(j))}$;
      append $\hat{\epsilon}_j$\;
      \Comment{slide only the season-$s(j)$ buffer}
    }
  }
}
\end{algorithm}
%------------------------------------------------------------

%------------------------------------------------------------
\subsection{Properties of \texttt{EnbPI-S}}
\label{sec:algorithm:properties}
%------------------------------------------------------------

The computational cost of \texttt{EnbPI-S} is negligibly different
from that of EnbPI.  The training phase is identical; in the
prediction phase, the only change is that the line
search~\eqref{eq:betahat_strat} and the quantile evaluations are
performed on a buffer of size $T_{\sstar} \approx T/\Sper$ rather
than $T$, so that for $\Sper = 12$ and the typical batch size
$s_0 = 1$ each step becomes cheaper by a factor of $\Sper$.  The
total memory required for the $\Sper$ buffers is still $T$ residuals.

No data splitting is needed either.  As in EnbPI, all $T$ training
observations are used to fit the $B$ bootstrap models, and all $T$
LOO residuals are used to populate the $\Sper$ stratified buffers.
This stands in contrast to Mondrian CP for i.i.d.\ data
\citep{vovk2005mondrian,dewolf2025conditional}, where the calibration
set must be divided among the $\Sper$ classes, which reduces the
effective calibration size to $T/\Sper$ per class.  In
\texttt{EnbPI-S} the training set is shared across seasons through
the LOO ensemble, and only the residual calibration is stratified.

The asymmetry of the interval is also season-adaptive.  The line
search~\eqref{eq:betahat_strat} selects an \emph{optimal asymmetric
quantile pair} $(\hat{\beta}^{(\sstar)},
1-\alpha+\hat{\beta}^{(\sstar)})$ tailored to the empirical
distribution of season-$\sstar$ residuals.  For seasons in which
forecast errors are right-skewed, as happens in January for the IPCA,
where large positive surprises dominate, the optimiser selects
$\hat{\beta}^{(\sstar)} < \alpha/2$ and produces a left-asymmetric
interval, adapting to distributional shape and not only to scale.
Section~\ref{sec:algorithm:defects} shows that this adaptation is not
free once the buffer is short.

Finally, the update rule~\eqref{eq:slide_strat} ensures that each
season-$s$ buffer always contains the $T_s$ most recent residuals
from season $s$.  \texttt{EnbPI-S} therefore tracks potential slow drift in
$\sigma_s$ over time, such as structural changes in seasonal
volatility patterns, while remaining robust to contamination from
other seasons.

%------------------------------------------------------------
\subsection{Normalisation as the Other Endpoint}
\label{sec:algorithm:normalised}
%------------------------------------------------------------

Stratification spends the seasonal information generously: it estimates
$\Sper$ complete distributions, one per season, from $\Ts$ observations
each.  Under Assumption~\ref{ass:seasonal} that is more than the problem
requires, since the seasons then share a shape and differ only in scale,
so $\Sper$ numbers suffice.  Let $\hat{\sigma}_s$ be a scale estimate
computed from the season-$s$ residuals, and let
\begin{equation}
  \label{eq:normalised_buffer}
  \hat{z}_t \;=\; \frac{\hat{\epsilon}_t}{\hat{\sigma}_{s(t)}},
  \qquad t = 1,\ldots,T,
\end{equation}
be the pooled buffer of standardised residuals.  The normalised scheme,
which we call \texttt{EnbPI-N}, runs the line search on
$\hat{\boldsymbol z}$ and multiplies the resulting offsets by
$\hat{\sigma}_{\sstar}$:
\begin{equation}
  \label{eq:enbpin_interval}
  \hat{C}_t^{\,\alpha,\textup{norm}}
  =
  \Bigl[
    \hat{f}_{-t}^{\,\phi}(x_t)
    + \hat{\sigma}_{\sstar}\Quant{\hat{\beta}}{\hat{\boldsymbol z}},\;
    \hat{f}_{-t}^{\,\phi}(x_t)
    + \hat{\sigma}_{\sstar}\Quant{1-\alpha+\hat{\beta}}{\hat{\boldsymbol z}}
  \Bigr].
\end{equation}
All $T$ residuals inform the shape and only $\Ts$ inform the scale, which
is the whole of the difference between the two schemes.  We use the median
absolute deviation divided by $\Phi^{-1}(0.75)$ for
$\hat{\sigma}_s$, since the residual distributions in our applications are
heavy-tailed, and freeze the scales at the end of the training phase.

If Assumption~\ref{ass:seasonal} holds, \texttt{EnbPI-N} removes the bias
of Proposition~\ref{prop:pooled_gap} just as stratification does, and it
does so from a longer buffer.  If it fails, that is, if seasons differ in
shape and not only in scale, the standardised residuals are still not
identically distributed, and normalisation leaves a discrepancy behind.
Write
\begin{equation}
  \label{eq:shape_discrepancy}
  \Delta_{\sstar}(p)
  \;:=\;
  \sigma_{\sstar}G_z^{-1}(p) - F_{\sstar}^{-1}(p),
\end{equation}
where $G_z$ is the CDF of the standardised errors pooled across seasons.
Under Assumption~\ref{ass:seasonal}, $\Delta_{\sstar}(p) = 0$ for every
season and every level; in general it does not vanish, and it is the
quantity that decides which of the two schemes is appropriate.

%------------------------------------------------------------
\subsection{One Family, Two Endpoints}
\label{sec:algorithm:family}
%------------------------------------------------------------

The two schemes estimate the same object, the season-$\sstar$ residual
quantile at level $p$, and nothing forces a choice between the estimates.
Write $\hat{\theta}^{S}(p)$ for the stratified estimate, the empirical
$p$-quantile of $\hat{\boldsymbol\epsilon}^{(\sstar)}$, and
$\hat{\theta}^{N}(p) = \hat{\sigma}_{\sstar}\Quant{p}{\hat{\boldsymbol z}}$
for the normalised one, and define, for $\lambda \in [0,1]$,
\begin{equation}
  \label{eq:family}
  \hat{\theta}^{\lambda}(p)
  \;=\;
  \lambda\,\hat{\theta}^{S}(p)
  \;+\;
  (1-\lambda)\,\hat{\theta}^{N}(p).
\end{equation}
The interval is built from $\hat{\theta}^{\lambda}$ exactly as before, with
a single line search run on the combined quantile function so that the
asymmetry is selected once rather than inherited from either scheme.
Setting $\lambda = 1$ returns \texttt{EnbPI-S} and $\lambda = 0$ returns
\texttt{EnbPI-N}, so the family contains both, and the question of which
scheme to use becomes the question of where in $[0,1]$ to sit.

Composing the two corrections rather than interpolating them would
achieve nothing: standardising the residuals of a season by a constant and
then multiplying the quantiles of that same season back by it returns the
stratified interval unchanged.  Only a genuine interpolation moves between
the two.

%------------------------------------------------------------
\subsection{Two Defects That Are Not About Seasons}
\label{sec:algorithm:defects}
%------------------------------------------------------------

Before the family can be assessed, two properties of the calibration step
have to be separated from the seasonal question, because both punish short
buffers and the stratified scheme is the one with short buffers.

The first is the inward bias of extreme empirical quantiles.  The
empirical $p$-quantile of $\Ts$ residuals cannot fall outside the range of
those residuals, so for $p$ near one it is biased downwards and for $p$
near zero upwards.  At $\Ts = 20$ and $\alpha = 0.10$ the bias is not
small: Section~\ref{sec:simulations:e6constr} measures it at $+0.28$ and $-0.35$
at the two ends in one of our designs, comparable in magnitude to the
seasonal discrepancy the calibration is trying to capture.  The standard
remedy is the finite-sample conformal endpoint, and it is worth stating in
the form in which it is meant to be used, as an \emph{order statistic}
rather than as a quantile level: for a buffer of size $n$ sorted as
$\hat{\epsilon}_{(1)} \leq \cdots \leq \hat{\epsilon}_{(n)}$, the upper
endpoint is $\hat{\epsilon}_{(k_{\text{hi}})}$ with
$k_{\text{hi}} = \lceil (n+1)(1-\alpha/2)\rceil$ and the lower endpoint is
$\hat{\epsilon}_{(k_{\text{lo}})}$ with
$k_{\text{lo}} = \lfloor (n+1)\alpha/2 \rfloor$.  Handing the ratio $k/n$
to an interpolating quantile function is a different object, and it undoes
most of the correction at the lower end, where $k$ is one or two.  The two
indices exist only for $n \geq 2/\alpha - 1$ and fall strictly inside the
sample only for $n \geq 4/\alpha - 1$;
Section~\ref{sec:simulations:e6constr} works out what happens between the two.

The second is a selection effect in the line search itself.  The optimiser
picks the narrowest of $m$ candidate quantile pairs, all computed from the
same residuals, and the narrowest of many noisy candidates is narrower than
it should be.  The bias grows with $m$ and shrinks with the buffer length,
so it is invisible on a pooled buffer of $T - p$ residuals and material on
a season buffer of $\Ts$.  Section~\ref{sec:simulations:e6constr} measures between four
and six percentage points of coverage at $\Ts$ between $20$ and $50$,
against under half a percentage point at $n = 500$.  This is a property of the
original algorithm, not of its seasonal variant, but it lands on the
seasonal variant.

Neither defect has anything to do with the seasonal bias of
Proposition~\ref{prop:pooled_gap}, and neither is repaired by choosing
$\lambda$ well.  Both are repaired directly, by taking the conformal order
statistics and by fixing $\beta = \alpha/2$ when the buffer is short, and
the experiments of Section~\ref{sec:simulations} report the family with
both repairs in place.  Experiment E6 shows that the repairs matter more
than the choice of $\lambda$: with them, a season-specific interval
attains its nominal level at $\Ts = 20$, which is the sample size at which
the literature, and the earlier sections of this paper, report
stratification as too expensive to use.

%------------------------------------------------------------
\subsection{Relation to Existing Methods}
\label{sec:algorithm:relation}
%------------------------------------------------------------

According to \citet{vovk2005mondrian}, Mondrian conformal prediction
stratifies the calibration set of an inductive conformal predictor by
a categorical label $\kappa(x, y)$, which guarantees that the
resulting intervals are conditionally valid within each category
under exchangeability.  \texttt{EnbPI-S} can be viewed as the
non-exchangeable time-series analogue of Mondrian CP with
$\kappa(x_t, y_t) = s(t)$, the calendar season.  Unlike Mondrian CP,
it operates without data splitting, since the bootstrap training uses
all observations, and provides coverage guarantees under temporal
dependence.

The \emph{Sequential Predictive Conformal Inference} algorithm of
\citet{xu2023sequential} constructs prediction intervals by
estimating the conditional quantile of the next residual through a
quantile random forest trained on past residuals; SPCI exploits
serial correlation in residuals to produce narrower intervals
whenever residuals are autocorrelated.  \texttt{EnbPI-S} and SPCI
address complementary sources of miscoverage, since SPCI corrects for
residual autocorrelation while \texttt{EnbPI-S} corrects for seasonal
heteroskedasticity, so that the two approaches are compatible and can
in principle be combined (see Section~\ref{sec:conclusion}).

\citet{dewolf2025conditional} show that, for i.i.d.\ data from a
location-scale family, normalised CP, which divides each residual by
an estimate of $\sigma(x)$, achieves conditional validity without
Mondrian stratification.  That result has a direct seasonal
counterpart: under Assumption~\ref{ass:seasonal} the standardised
errors $\epsilon_t/\sigma_{s(t)}$ are identically distributed, so a
scheme that divides each residual by an estimate of
$\sigma_{s(t)}$ and calibrates on the pooled standardised buffer also
removes the bias of Proposition~\ref{prop:pooled_gap}, and does so
while spending only $\Sper$ parameters on the seasonal structure.
Experiment E7 in Section~\ref{sec:simulations:e7schemes} implements that
competitor, called \texttt{EnbPI-N}, and finds it the more efficient of
the two whenever the scale-mixture restriction holds, delivering a
comparable coverage level from a narrower interval.  The case for
stratification is therefore not that it dominates normalisation, but
that it does not depend on the restriction:
Assumption~\ref{ass:seasonal} is a modelling choice, and when seasons
differ in the shape of the error distribution rather than in its scale,
a scale correction has nothing to correct while stratification still
equalises coverage.  Experiment E6 quantifies that case as well, and
Section~\ref{sec:empirical:choice} tests which of the two
configurations the Brazilian data resemble.
